设三场满足命题 \ref{prop:evolution} 的条件, 普通连续谱上 $s_{11}s_{22}\ne0$, 离散谱有限且简单. 记 $s_{22}$ 在 $D_+$ 中的全部零点为 $\zeta_1,\ldots,\zeta_m$, $m=2N$, 范数常数及 $C_j=c_je^{2\ii\theta(\zeta_j)}$ 按第 \ref{ch:discrete} 节定义. 将 $\Sigma$ 定向为 $D_+$ 位于左侧: 实轴的外部两段向右, 内部区间向左, 上半圆顺时针, 下半圆逆时针. 下标 $+$、$-$ 分别表示左、右非切向边界值. 普通轮廓指删去 $0,\pm q_0,\pm\ii q_0,z_a^\pm$ 后的各弧.

\begin{theorem}[非零背景的矩阵 RHP]\label{thm:RHP}\label{prop:jump}
由 Jost 解析列定义
\begin{equation}\label{eq:Mdef}
M(z)=\begin{cases}
M_+(z)=\left(\mu_{-,1},\dfrac{\mu_{+,2}}{s_{22}}\right),&z\in D_+,\\[6pt]
M_-(z)=\left(\dfrac{\mu_{+,1}}{s_{11}},\mu_{-,2}\right),&z\in D_-.
\end{cases}
\end{equation}
则 $M$ 在两域内亚纯, $D_+$ 中仅第二列可在 $\zeta_j$ 处有简单极点, $D_-$ 中仅第一列可在 $\conj{\zeta_j}$ 处有简单极点. 普通轮廓上的跳跃为
\begin{equation}\label{eq:jump}
M_+=M_-J,\qquad
J=\begin{pmatrix}
1&\widetilde\rho\,e^{2\ii\theta}\\
-\rho\,e^{-2\ii\theta}&1-\rho\widetilde\rho
\end{pmatrix},\qquad \det J=1.
\end{equation}
各简单极点的留数满足 \eqref{eq:resplus}--\eqref{eq:resminus}. 行列式、两处归一化与约化为
\begin{align}
\det M(z)&=\gamma(z),\label{eq:detM}\\
M(z)&=I+\dfrac{M^{(1)}}{z}+O(z^{-2}),\qquad z\to\infty,\label{eq:Minfty}\\
M(z)&=-\dfrac{\ii q_0}{z}\sigma_1+O(1),\qquad z\to0,\label{eq:Mzero}\\
M(z)&=\sigma_2\conj{M(\conj{z})}\sigma_2,\qquad
M(h(z))=M(z)\dfrac{\ii z}{q_0}\sigma_1.\label{eq:Msym}
\end{align}
\end{theorem}
\begin{proof}
解析性与列极点由第 \ref{ch:direct}、\ref{ch:scattering}、\ref{ch:discrete} 节的结果得到. 去除散射关系中的指数, 有
\[
\mu_{+,1}=s_{11}\mu_{-,1}+s_{21}e^{-2\ii\theta}\mu_{-,2},\qquad
\mu_{+,2}=s_{12}e^{2\ii\theta}\mu_{-,1}+s_{22}\mu_{-,2}.
\]
第一式给出 $M_{+,1}=M_{-,1}-\rho e^{-2\ii\theta}M_{-,2}$, 代入第二式得到
\[
M_{+,2}=\widetilde\rho e^{2\ii\theta}M_{-,1}+(1-\rho\widetilde\rho)M_{-,2},
\]
即 \eqref{eq:jump}. Wronskian 表达式给出
\[
\det M_+=\dfrac{\det(\mu_{-,1},\mu_{+,2})}{s_{22}}=\gamma,
\]
另一域同理. 第 \ref{ch:asymptotics} 节的展开确定两处归一化. 共轭约化交换两列和两域; 换叶约化将 $\mu_{-,1}$ 与 $\mu_{-,2}$、$\dfrac{\mu_{+,2}}{s_{22}}$ 与 $\dfrac{\mu_{+,1}}{s_{11}}$ 配对, 因而得到 \eqref{eq:Msym}.
\end{proof}

\subsection{留数条件与重构公式}
记 $e_1,e_2$ 为标准基, 置
\begin{equation}\label{eq:Gdefinition}
G=I-J^{-1}=\begin{pmatrix}
\rho\widetilde\rho&\widetilde\rho e^{2\ii\theta}\\
-\rho e^{-2\ii\theta}&0
\end{pmatrix},\qquad M_+-M_-=M_+G.
\end{equation}
以下表示按照先消去奇性、再应用 Plemelj 公式的顺序建立\cite{ZhangFan2020,Fan2022}. 它适用于满足所列条件的 RHP 解.

\begin{proposition}[含反射的留数积分表示]\label{prop:cauchy-representation}
设定理 \ref{thm:RHP} 的 RHP 有解. 假设减去离散主部及 $z=0$ 主部后, 矩阵在各有限特殊点有有界、相容的扇区边界值; 跳跃密度 $W=M_+G$ 在普通弧上局部 Hölder 连续, 且 $W,sW\in L^1(\Sigma)$. 删去特殊点的小圆弧时, 假设其上的正则化矩阵一致有界. 则
\begin{equation}\label{eq:general-cauchy}
\begin{aligned}
M(z)={}&I-\dfrac{\ii q_0}{z}\sigma_1
+\sum\limits_{j=1}^{m}\left(
\dfrac{C_j\mathbf v_je_2^T}{z-\zeta_j}
-\dfrac{\conj{C_j}\mathbf u_je_1^T}{z-\conj{\zeta_j}}\right)\\
&+\dfrac{1}{2\pi\ii}\int\limits_\Sigma\dfrac{M_+(s)G(s)}{s-z}\,\dd s,
\qquad z\in\mathbb C\mathbin{\backslash}\Sigma,
\end{aligned}
\end{equation}
其中 $\mathbf v_j=M_1(\zeta_j)$、$\mathbf u_j=M_2(\conj{\zeta_j})$ 是各点的正则列. 它们满足积分--代数关系
\begin{align}
\mathbf v_j={}&e_1-\dfrac{\ii q_0}{\zeta_j}e_2
-\sum\limits_{k=1}^{m}\dfrac{\conj{C_k}\mathbf u_k}{\zeta_j-\conj{\zeta_k}}
+\dfrac{1}{2\pi\ii}\int\limits_\Sigma\dfrac{W_1(s)}{s-\zeta_j}\,\dd s,
\label{eq:general-finite1}\\
\mathbf u_j={}&e_2-\dfrac{\ii q_0}{\conj{\zeta_j}}e_1
+\sum\limits_{k=1}^{m}\dfrac{C_k\mathbf v_k}{\conj{\zeta_j}-\zeta_k}
+\dfrac{1}{2\pi\ii}\int\limits_\Sigma\dfrac{W_2(s)}{s-\conj{\zeta_j}}\,\dd s.
\label{eq:general-finite2}
\end{align}
按 $q=\ii M^{(1)}_{12}$ 定义重构光场, 则
\begin{equation}\label{eq:qgeneral}
q=q_0+\ii\sum\limits_{j=1}^{m}C_j(\mathbf v_j)_1
-\dfrac{1}{2\pi}\int\limits_\Sigma W_{12}(s)\,\dd s.
\end{equation}
\end{proposition}
\begin{proof}
第 \ref{ch:discrete} 节的留数条件给出
\[
\Res_{z=\zeta_j}M=C_j\mathbf v_je_2^T,\qquad
\Res_{z=\conj{\zeta_j}}M=-\conj{C_j}\mathbf u_je_1^T.
\]
从 $M$ 减去 $I-\dfrac{\ii q_0}{z}\sigma_1$ 及上述各主部, 所得分片解析矩阵 $F_0$ 在无穷远为 $O(z^{-1})$, 而 $F_{0,+}-F_{0,-}=W$. 对删去特殊点小圆盘并截断于大圆的两域分别使用 Cauchy 公式. 左侧为正向边界, 将两域的公共边界合并得到 $W$; 局部一致有界性保证小圆贡献消失, 大圆上的 $O(z^{-1})$ 使其贡献消失. 因而
\[
F_0(z)=\dfrac{1}{2\pi\ii}\int\limits_\Sigma\dfrac{W(s)}{s-z}\,\dd s,
\]
得到 \eqref{eq:general-cauchy}. 正号由 \eqref{eq:Gdefinition} 与轮廓定向共同确定.

在第一列取 $z=\zeta_j$, 在第二列取 $z=\conj{\zeta_j}$, 各自本列在该点没有离散极点, 得 \eqref{eq:general-finite1}--\eqref{eq:general-finite2}. 沿避开轮廓的闭角域令 $z\to\infty$, 使用
\[
\dfrac{1}{s-z}=-\dfrac{1}{z}+\dfrac{s}{z(s-z)},\qquad
\left|\dfrac{s}{z(s-z)}\right|\le\dfrac{C|s|}{|z|^2},
\]
及 $sW\in L^1$ 控制余项, 得
\begin{equation}\label{eq:M1-general}
M^{(1)}=-\ii q_0\sigma_1
+\sum\limits_{j=1}^{m}(C_j\mathbf v_je_2^T-\conj{C_j}\mathbf u_je_1^T)
-\dfrac{1}{2\pi\ii}\int\limits_\Sigma W(s)\,\dd s.
\end{equation}
读取 $(1,2)$ 元的系数并乘以 $\ii$, 即得 \eqref{eq:qgeneral}. 对直接散射构造的 $M$, 此定义与原光场一致, 见下述重构命题.
\end{proof}

\begin{proposition}[光场与介质重构]\label{prop:reconstruction}
直接散射构造的矩阵 $M$ 满足
\begin{equation}\label{eq:reconq}
Q(q)=\dfrac{\ii}{2}[\sg,M^{(1)}],\qquad q=\ii M^{(1)}_{12}.
\end{equation}
若满足命题 \ref{prop:resonance} 的共振相容条件, 则 $M$ 在两个共振原像的空间归一化值可逆, 且
\begin{equation}\label{eq:RfromM}
R=M(x,t,z_a)\dfrac{\eta_0\lambda(z_a)}{a}\sg M(x,t,z_a)^{-1},
\qquad z_a=z_a^\pm.
\end{equation}
两式给出同一个介质矩阵. 原物理场由
\begin{equation}\label{eq:original-fields}
\widetilde q=e^{-\ii\Omega t}q,\qquad
\widetilde p=-e^{-\ii\Omega t}R_{12},\qquad \eta=R_{11}
\end{equation}
恢复.
\end{proposition}
\begin{proof}
由 \eqref{eq:mu-highenergy} 与 \eqref{eq:transmission-highenergy}, 在 $D_+$ 中第二列除以 $s_{22}$ 后, 其对角系数从 $-\ii J_+$ 变为 $-\ii(J_++J_{\rm tot})=-\ii J_-$. 另一域相同, 故
\[
M^{(1)}=-\ii B(q)+\ii J_-\sg,
\]
得到 \eqref{eq:reconq}. 这一关系也可由空间方程
\begin{equation}\label{eq:Mspace}
M_x=-\ii k\sg M+QM+\ii\lambda M\sg
\end{equation}
比较无穷远的常数项得到: 第 \ref{ch:asymptotics} 节的积分估计还给出 $M_x=O(z^{-1})$, 而 $k,\lambda\sim\dfrac{z}{2}$.

在共振点, $S_{\rm sp}$ 对角, 两种拼接均为 $M=\mu_-$. 命题 \ref{prop:resonance} 的证明给出
\[
\chi_-^{-1}R\chi_-=\dfrac{\eta_0\lambda(z_a)}{a}\sg.
\]
空间指数与 $\sg$ 对易, 因而得到 \eqref{eq:RfromM}; 这里使用空间归一化值, 不对奇异的传播指数求值. $\gamma(z_a)>0$ 保证可逆. 换叶时 $\lambda$ 变号, 同时 $\sigma_1\sg\sigma_1=-\sg$, 两次变号相消, 故两原像的结果一致.
\end{proof}

\subsection{迹公式与相位条件}
相位条件由 $z=0$ 处的归一化确定. 透射系数同时包含离散零点和连续谱反射的贡献. 以下标量分解采用非零边界逆散射中的 Cauchy 积分方法\cite{Fan2022,ZhangFan2020}, 积分轮廓取 HMB 的实轴与圆周.

\begin{proposition}[迹公式与同相约束]\label{prop:theta}
设透射系数在普通轮廓上非零, 只有上述有限简单离散零点. 取
$\ell(s)=\log(1-\rho(s)\widetilde\rho(s))$ 的零指标分支, 设其在普通弧上局部 Hölder 连续, 在 $0,\infty$ 的极限为零, 且 $\ell,\dfrac{\ell(s)}{s}\in L^1(\Sigma)$. 假设剔除零点后两域中的透射函数具有与乘积关系相容的单值对数, 这些对数与下述 Cauchy 积分在各有限特殊点的边界值有界并局部匹配. 则
\begin{equation}\label{eq:trace}
s_{22}(z)=\prod_{j=1}^{m}\dfrac{z-\zeta_j}{z-\conj{\zeta_j}}
\exp\left\{-\dfrac{1}{2\pi\ii}\int\limits_\Sigma\dfrac{\ell(s)}{s-z}\,\dd s\right\},\qquad z\in D_+.
\end{equation}
$s_{11}$ 由共轭约化确定. 相同左右背景要求
\begin{equation}\label{eq:theta-condition}
\prod_{j=1}^{m}\dfrac{\zeta_j}{\conj{\zeta_j}}
\exp\left\{-\dfrac{1}{2\pi\ii}\int\limits_\Sigma\dfrac{\ell(s)}{s}\,\dd s\right\}=1.
\end{equation}
特别地, 无反射时 $s_{22}$ 等于 \eqref{eq:trace} 中的有理乘积, 且 $\displaystyle\prod_{j=1}^{m}\dfrac{\zeta_j}{\conj{\zeta_j}}=1$.
\end{proposition}
\begin{proof}
记 $B_0(z)=\displaystyle\prod_{j=1}^{m}\dfrac{z-\zeta_j}{z-\conj{\zeta_j}}$, 定义
$\beta_+=\dfrac{s_{22}}{B_0}$、$\beta_-=s_{11}B_0$. 它们在相应域解析且无零点, 无穷远归一化为 1. 由 $\det S=1$,
\[
\beta_+\beta_-=s_{11}s_{22}=\dfrac{1}{1-\rho\widetilde\rho}.
\]
相容的对数分支满足 $\log\beta_++\log\beta_-=-\ell$, 因而分片解析函数 $f_+=\log\beta_+$、$f_-=-\log\beta_-$ 的跳跃为 $f_+-f_-=-\ell$. Plemelj 公式及归一化给出
\[
f(z)=-\dfrac{1}{2\pi\ii}\int\limits_\Sigma\dfrac{\ell(s)}{s-z}\,\dd s.
\]
差函数在特殊点可去且在无穷远为零, 因而没有附加整函数项. 代回 $\beta_+$ 即得 \eqref{eq:trace}. 为取零点极限, 沿 $D_+$ 内的 $z=-\ii\varepsilon$, $\varepsilon\downarrow0$. 实轴上 $\left|\dfrac{s}{s+\ii\varepsilon}\right|\le1$, 圆周与零点分离, 故 $\dfrac{\ell(s)}{s-z}$ 受 $C\left|\dfrac{\ell(s)}{s}\right|$ 控制. 支配收敛结合 $s_{22}(0)=1$ 给出 \eqref{eq:theta-condition}.
\end{proof}

每个完整四点轨道取外部上半平面的代表 $\zeta_n$, 则它的两个 $D_+$ 零点对无反射相位乘积的贡献为 $\left(\dfrac{\zeta_n}{\conj{\zeta_n}}\right)^2$. 因而同相约束化为
\[
4\sum\limits_{n=1}^{N}\arg\zeta_n=0\pmod{2\pi}.
\]
纯虚代表逐轨道满足这一条件.

\subsection{\texorpdfstring{$N$}{N} 个完整简单四点轨道的无反射 RHP 求解}
当 $\rho=\widetilde\rho=0$ 时, 跳跃矩阵为 $I$. 在非极点轮廓处取可去条件. 在分支点 $z_b=\pm\ii q_0$ 处, 矩阵 $M$ 正则, 而其行列式 $\gamma(z)=1+\dfrac{q_0^2}{z^2}$ 具有一阶零点, 即 $\gamma(z_b)=0$、$\gamma'(z_b)\ne0$.

固定 $N$ 个两两不交的完整简单四点轨道. 按 \eqref{eq:discrete-indexing} 选取外部上半平面代表 $\zeta_1,\ldots,\zeta_N$, 并置
\begin{equation}\label{eq:N-orbit-data}
\begin{aligned}
m&=2N,\qquad \zeta_{N+n}=-\dfrac{q_0^2}{\conj{\zeta_n}},\qquad n=1,\ldots,N,\\
c_{N+n}&=-\dfrac{q_0^2}{(\conj{\zeta_n})^2}\conj{c_n},\qquad
C_j=c_je^{2\ii\theta(\zeta_j)},\qquad j=1,\ldots,m.
\end{aligned}
\end{equation}
所有 $4N$ 个谱点互异且位于 $\Sigma$ 外, 各 $c_n\ne0$. 同相背景要求
\begin{equation}\label{eq:N-orbit-phase}
\prod\limits_{j=1}^{m}\dfrac{\zeta_j}{\conj{\zeta_j}}=1.
\end{equation}
下文保留全部 $m$ 个 $D_+$ 点及其共轭点, 令 $\mathbf v_j=M_1(\zeta_j)$、$\mathbf u_j=M_2(\conj{\zeta_j})$.

\begin{proposition}[留数系统]\label{prop:finite}
除 $0$ 与指定离散极点外没有有限奇点的无反射矩阵解具有表示
\begin{equation}\label{eq:partialfraction}
M(z)=I-\dfrac{\ii q_0}{z}\sigma_1+\sum\limits_{j=1}^{m}
\left(\dfrac{C_j\mathbf v_je_2^T}{z-\zeta_j}
-\dfrac{\conj{C_j}\mathbf u_je_1^T}{z-\conj{\zeta_j}}\right),
\end{equation}
其正则列满足
\begin{align}
\mathbf v_j&=e_1-\dfrac{\ii q_0}{\zeta_j}e_2
-\sum\limits_{k=1}^{m}\dfrac{\conj{C_k}\mathbf u_k}{\zeta_j-\conj{\zeta_k}},\label{eq:finite1}\\
\mathbf u_j&=e_2-\dfrac{\ii q_0}{\conj{\zeta_j}}e_1
+\sum\limits_{k=1}^{m}\dfrac{C_k\mathbf v_k}{\conj{\zeta_j}-\zeta_k}.\label{eq:finite2}
\end{align}
反之, 任一系统解通过 \eqref{eq:partialfraction} 给出具有指定留数和两处归一化的矩阵. 记
\[
A_{jk}=\dfrac{\conj{C_k}}{\zeta_j-\conj{\zeta_k}},\qquad
B_{jk}=\dfrac{C_k}{\conj{\zeta_j}-\zeta_k},\qquad j,k=1,\ldots,m,
\]
并记 $f_j=e_1-\dfrac{\ii q_0}{\zeta_j}e_2$、$g_j=e_2-\dfrac{\ii q_0}{\conj{\zeta_j}}e_1$. 将 $\mathbf v_j^T,\mathbf u_j^T,f_j^T,g_j^T$ 分别作为 $m\times2$ 矩阵 $\boldsymbol v,\boldsymbol u,\boldsymbol f,\boldsymbol g$ 的第 $j$ 行. 则留数系统为
\begin{equation}\label{eq:N-block-system}
K_N\begin{pmatrix}\boldsymbol v\\\boldsymbol u\end{pmatrix}
=\begin{pmatrix}\boldsymbol f\\\boldsymbol g\end{pmatrix},\qquad
K_N=\begin{pmatrix}I_m&A\\-B&I_m\end{pmatrix}.
\end{equation}
置 $L_N=I_m+AB$. 若 $\det L_N\ne0$, 则系统的唯一解为
\begin{equation}\label{eq:finite-schur}
\boldsymbol v=L_N^{-1}(\boldsymbol f-A\boldsymbol g),\qquad
\boldsymbol u=\boldsymbol g+B\boldsymbol v.
\end{equation}
相应光场为
\begin{equation}\label{eq:qfinite}
q=q_0+\ii\sum\limits_{j=1}^{m}C_j(\mathbf v_j)_1.
\end{equation}
\end{proposition}
\begin{proof}
减去 $z=0$ 处主部及全部离散主部后, 剩余矩阵在有限平面解析, 且在无穷远趋于 $I$. Liouville 定理给出 \eqref{eq:partialfraction}, 这也正是 \eqref{eq:general-cauchy} 中 $G=0$ 的情形. 在第一列取 $z=\zeta_j$, 在第二列取 $z=\conj{\zeta_j}$, 即得 \eqref{eq:finite1}--\eqref{eq:finite2}. 反向代入时, 这些取值等式保证留数系数恰为对应正则列. 写成
\[
\boldsymbol v=\boldsymbol f-A\boldsymbol u,\qquad
\boldsymbol u=\boldsymbol g+B\boldsymbol v,
\]
对系数矩阵作分块消元,
\[
\begin{pmatrix}I_m&-A\\0&I_m\end{pmatrix}
\begin{pmatrix}I_m&A\\-B&I_m\end{pmatrix}
=\begin{pmatrix}L_N&0\\-B&I_m\end{pmatrix},
\qquad \det K_N=\det L_N.
\]
当此行列式非零时, 第一行解出 $\boldsymbol v$, 第二行给出 $\boldsymbol u$, 得到 \eqref{eq:finite-schur}. 最后读取 \eqref{eq:partialfraction} 的 $(1,2)$ 元的 $z^{-1}$ 系数, 得到 \eqref{eq:qfinite}.
\end{proof}

\begin{theorem}[$N$ 轨道矩阵解与行列式重构]\label{thm:N-orbit-RHP}
设数据满足 \eqref{eq:N-orbit-data}--\eqref{eq:N-orbit-phase}. 在任一满足 $\det L_N(x,t)\ne0$ 的实点 $(x,t)$, 式 \eqref{eq:finite-schur} 和 \eqref{eq:partialfraction} 给出无反射 RHP 的唯一解. 它满足 \eqref{eq:Msym} 及 $\det M=\gamma$, 在分支点解析且秩为一. 令
\begin{equation}\label{eq:N-rhs}
\begin{aligned}
&\boldsymbol C=(C_1,\ldots,C_m)^T,\qquad
\boldsymbol r=(\boldsymbol f-A\boldsymbol g)e_1,\\
&r_j=1+\ii q_0\sum\limits_{k=1}^{m}
\dfrac{\conj{C_k}}{(\zeta_j-\conj{\zeta_k})\conj{\zeta_k}},\qquad j=1,\ldots,m,
\end{aligned}
\end{equation}
则规范光场有显式矩阵式与行列式表示
\begin{equation}\label{eq:N-determinant-q}
q=q_0+\ii\boldsymbol C^T L_N^{-1}\boldsymbol r
=\dfrac{\det\begin{pmatrix}L_N&\boldsymbol r\\-\ii\boldsymbol C^T&q_0\end{pmatrix}}
{\det L_N}.
\end{equation}
介质矩阵由同一有理矩阵在共振原像处重构:
\begin{equation}\label{eq:N-medium}
R=M(z_a^\pm)\dfrac{\eta_0\lambda(z_a^\pm)}{a}\sigma_3 M(z_a^\pm)^{-1}.
\end{equation}
两个符号给出相同结果; 原三场按 \eqref{eq:original-fields} 恢复.
\end{theorem}
\begin{proof}
命题 \ref{prop:finite} 给出有理矩阵的全部留数及归一化, 并由 $\det L_N\ne0$ 保证唯一性. 共轭变换 $\sigma_2\conj{M(\conj{z})}\sigma_2$ 交换两组极点, 满足同一留数系统和两处归一化. 换叶变换
\[
\widehat M(z)=M(h(z))\left(\dfrac{\ii z}{q_0}\sigma_1\right)^{-1}
\]
也满足这些条件. 具体地, 记同一轨道的配对指标为 $\widehat j$, 即 $\zeta_{\widehat j}=h(\conj{\zeta_j})$. 在 $z=\zeta_j$ 处, $h(z)$ 趋于 $\conj{\zeta_{\widehat j}}$, 由 $h'(\zeta_j)=\dfrac{q_0^2}{\zeta_j^2}$ 和 \eqref{eq:Cpair},
\[
\Res_{z=\zeta_j}\widehat M_2(z)
=\dfrac{\ii q_0\conj{C_{\widehat j}}}{\zeta_jh'(\zeta_j)}\mathbf u_{\widehat j}
=C_j\left(-\dfrac{\ii q_0}{\zeta_j}\mathbf u_{\widehat j}\right)
=C_j\widehat M_1(\zeta_j).
\]
共轭点处同理, 得到带负号的第一列留数条件. 在 $0,\infty$ 处, 换叶交换两处归一化, 仍给出规定主项. 由唯一性得到两项约化 \eqref{eq:Msym}.

在每个离散极点, 奇异列的留数为正则列的标量倍数, 因而 $\det M$ 的留数为零. 又由 $0,\infty$ 的归一化,
\[
\det M(z)=1+\dfrac{d}{z}+\dfrac{q_0^2}{z^2}
\]
对某个标量 $d$ 成立. 换叶约化的行列式给出
\[
1-\dfrac{dz}{q_0^2}+\dfrac{z^2}{q_0^2}
=\det M(h(z))=\dfrac{z^2}{q_0^2}\det M(z)
=\dfrac{z^2}{q_0^2}+\dfrac{dz}{q_0^2}+1,
\]
故 $d=0$, 即 $\det M=\gamma$. 分支点不在部分分式的极点集合中, 所以 $M$ 在该处解析. 若 $M(z_b)$ 为零矩阵, 则行列式至少有二阶零点, 与 $\gamma'(z_b)\ne0$ 矛盾; 因此其秩恰为一.

由 \eqref{eq:finite-schur}, 正则列第一分量排成的列向量为 $L_N^{-1}\boldsymbol r$. 代入 \eqref{eq:qfinite} 得到 \eqref{eq:N-determinant-q} 的第一式. 再对其分子作分块行列式计算,
\[
\det\begin{pmatrix}L_N&\boldsymbol r\\-\ii\boldsymbol C^T&q_0\end{pmatrix}
=\det L_N\left(q_0+\ii\boldsymbol C^T L_N^{-1}\boldsymbol r\right),
\]
得到第二式. 共振原像不在离散谱中, 且 $\gamma(z_a^\pm)>0$, 所以逆矩阵存在. 两个原像互为换叶, $\lambda(h(z))=-\lambda(z)$, 而 $\sigma_1\sigma_3\sigma_1=-\sigma_3$; 代入换叶约化即得 \eqref{eq:N-medium} 与谱叶无关.
\end{proof}

上述公式将 $N$ 轨道问题化为 $2N$ 阶矩阵 $L_N$ 的求逆. 公式适用于 $\det L_N\ne0$ 的区域; 下节对一个纯虚轨道进一步显式消元, 并证明该行列式在全体实 $x,t$ 上严格为正.

\subsection{单轨道解及 HMB 方程验证}
取一个纯虚离散轨道
\begin{equation}\label{eq:imaginarypoint}
k_1=\ii b,\qquad b>q_0,\qquad
\kappa=\sqrt{b^2-q_0^2},\qquad W=b+\kappa,\qquad h_0=\dfrac{q_0}{W}\in(0,1).
\end{equation}
两个 $D_+$ 零点与透射函数为
\begin{equation}\label{eq:onezeros}
\zeta_1=\ii W,\qquad \zeta_2=-\ii q_0h_0=h(\conj{\zeta_1}),\qquad
s_{22}(z)=\dfrac{\lambda(z)-\ii\kappa}{\lambda(z)+\ii\kappa}.
\end{equation}
这正是无反射迹公式的两因子乘积, 其在 $0,\infty$ 均趋于 1. 在外部零点选择 $b_1=1$, 并由换叶约化确定另一常数:
\begin{equation}\label{eq:onec}
c_1=\dfrac{1}{s_{22}'(\zeta_1)}=\dfrac{4\ii\kappa}{1+h_0^2},\qquad
c_2=-h_0^2c_1.
\end{equation}
令
\begin{equation}\label{eq:orbit-phase}
s=\kappa[x-F(\ii b)t],\qquad T=e^{2s}.
\end{equation}
由 $\lambda(\zeta_1)=\lambda(\zeta_2)=\ii\kappa$, $k(\zeta_2)=-\ii b$ 及 $F(-\ii b)=\conj{F(\ii b)}$, 得 $C_1=c_1T$、$C_2=-h_0^2c_1\conj{T}$.

\begin{lemma}[一轨道系统的消元]\label{lem:direct-solve}
上述留数系统对任意实 $x,t$ 唯一可解. 记
\begin{equation}\label{eq:delta-pair}
\Delta_\pm=(1+h_0^2)(1+|T|^2)\pm2h_0(T+\conj{T}),
\end{equation}
则其解为
\begin{align}
\mathbf v_1&=\dfrac{1+h_0^2}{\Delta_-}
\begin{pmatrix}1-h_0\conj{T}\\\conj{T}-h_0\end{pmatrix},&
\mathbf u_1&=\dfrac{1+h_0^2}{\Delta_-}
\begin{pmatrix}h_0-T\\1-h_0T\end{pmatrix},\label{eq:direct-vu1}\\
\mathbf v_2&=\dfrac{1}{h_0}\mathbf u_1,&
\mathbf u_2&=-\dfrac{1}{h_0}\mathbf v_1.\label{eq:direct-vu2}
\end{align}
\end{lemma}
\begin{proof}
代入谱点和范数常数得到
\begin{equation}\label{eq:one-AB}
A=\dfrac{1}{1+h_0^2}
\begin{pmatrix}(h_0^2-1)\conj{T}&2h_0^2T\\-2\conj{T}&(h_0^2-1)T\end{pmatrix},\quad
B=\dfrac{1}{1+h_0^2}
\begin{pmatrix}(h_0^2-1)T&2h_0^2\conj{T}\\-2T&(h_0^2-1)\conj{T}\end{pmatrix}.
\end{equation}
对应右端为
\[
f_1=\begin{pmatrix}1\\-h_0\end{pmatrix},\quad
f_2=\begin{pmatrix}1\\h_0^{-1}\end{pmatrix},\quad
g_1=\begin{pmatrix}h_0\\1\end{pmatrix},\quad
g_2=\begin{pmatrix}-h_0^{-1}\\1\end{pmatrix}.
\]
两个向量分量分别满足同一个四阶标量系数系统,
\[
\begin{pmatrix}I_2&A\\-B&I_2\end{pmatrix}
\begin{pmatrix}\boldsymbol v\\\boldsymbol u\end{pmatrix}
=\begin{pmatrix}\boldsymbol f\\\boldsymbol g\end{pmatrix}.
\]
其 Schur 补行列式为
\begin{equation}\label{eq:one-schur-det}
\det(I_2+AB)=\dfrac{\Delta_-\Delta_+}{(1+h_0^2)^2}>0.
\end{equation}
事实上,
\[
\Delta_-=|T-h_0|^2+|1-h_0T|^2,\qquad
\Delta_+=|T+h_0|^2+|1+h_0T|^2,
\]
且 $\Delta_\pm\ge(1-h_0)^2(1+|T|^2)$. 令 $L=I_2+AB$, 用二阶逆矩阵公式计算 $L^{-1}(\boldsymbol f-A\boldsymbol g)$, 约去公因子 $\Delta_+$, 得
\[
\boldsymbol v=\dfrac{1+h_0^2}{\Delta_-}
\begin{pmatrix}1-h_0\conj{T}&\conj{T}-h_0\\
\dfrac{h_0-T}{h_0}&\dfrac{1-h_0T}{h_0}\end{pmatrix}.
\]
此处每一行是一个正则列的转置. 由 $\boldsymbol u=\boldsymbol g+B\boldsymbol v$ 得到其余两列, 即 \eqref{eq:direct-vu1}--\eqref{eq:direct-vu2}.
\end{proof}

将所得列向量代入部分分式表示. 其系数可用以下秩一正交投影写成紧凑形式:
\begin{equation}\label{eq:algebraic-P}
P=\dfrac{1}{\Delta_-}
\begin{pmatrix}
|T-h_0|^2&(T-h_0)(1-h_0\conj{T})\\
(1-h_0T)(\conj{T}-h_0)&|1-h_0T|^2
\end{pmatrix},\qquad H=\diag(\ii,-\ii).
\end{equation}
由 $\Delta_-=|T-h_0|^2+|1-h_0T|^2$, 有 $P^\dagger=P$, $P^2=P$, $\tr P=1$.

\begin{theorem}[完整四点轨道的矩阵解]\label{thm:oneRHP}\label{prop:unique-reflectionless}
在背景参数 $q_0>0$, $a\ne0$, $\eta_0\ne0$, $\alpha,\beta\in\mathbb R$ 和 $b>q_0$ 下, 数据 \eqref{eq:onezeros}--\eqref{eq:onec} 的无反射 RHP 有唯一解
\begin{equation}\label{eq:explicitM}
M(z)=H\bigl[(k(z)+\ii b)I-2\ii bP\bigr]Y(z)
\diag\left(\dfrac{1}{\lambda(z)+\ii\kappa},\dfrac{1}{\lambda(z)-\ii\kappa}\right)H^{-1}.
\end{equation}
它满足定理 \ref{thm:RHP} 的全部条件, 在分支点、轮廓交点与共振原像解析, 在分支点秩为一, 在共振原像可逆. 唯一性取于除 $0$ 和指定四点外没有有限奇点、非极点轮廓处可去并满足留数及两处归一化的亚纯矩阵类.
\end{theorem}
\begin{proof}
在部分分式的第二列合并分母,
$(z-\zeta_1)(z-\zeta_2)=2z(\lambda-\ii\kappa)$, 第一列同理. 利用
$2bh_0=q_0(1+h_0^2)$、$2\kappa h_0=q_0(1-h_0^2)$ 化简分子, 得到 \eqref{eq:explicitM}. 留数、两处归一化及恒等跳跃已由命题 \ref{prop:finite} 保证.

共轭矩阵 $\sigma_2\conj{M(\conj{z})}\sigma_2$ 满足同一留数系统, 其中交换两组极点产生 \eqref{eq:resminus} 的负号. 换叶矩阵 $M(h(z))C(z)^{-1}$ 也满足该系统: 留数变换的导数因子由 \eqref{eq:cpair} 抵消, 而 $\theta(h(\conj{\zeta}))=-\conj{\theta(\zeta)}$ 保证指数一致. 两处归一化相互交换后仍为规定主项. 引理 \ref{lem:direct-solve} 的唯一性遂给出 \eqref{eq:Msym}.

记
\begin{equation}\label{eq:Nfactor}
N(k)=(k+\ii b)I-2\ii bP.
\end{equation}
$P$ 的特征值为 $1,0$, 故
\[
\det N=(k-\ii b)(k+\ii b)=\lambda^2+\kappa^2.
\]
它抵消两个列分母的乘积, 得 $\det M=\det Y=\gamma$. 在分支点 $z=\pm\ii q_0$ 上, $\lambda=0$, 两个列分母均非零; $N$ 可逆而 $Y$ 秩一, 因此 $M$ 解析且秩一. 交点和共振原像的 $\lambda$ 为实数, 也没有列分母零点; 共振处 $\gamma>0$, 所以 $M$ 可逆. 显式表达式不再含有其他有限奇点.

任一属于所述函数类的解均具有 \eqref{eq:partialfraction} 的表示, 因而由同一留数系统确定. Schur 行列式 \eqref{eq:one-schur-det} 严格为正, 保证该解唯一.
\end{proof}

重构公式将同一矩阵的无穷远系数与共振值分别转化为光场和介质. 以下结论同时给出三个物理场及其 Lax 相容性.

\begin{theorem}[单轨道三场解]\label{thm:direct-HMB}\label{prop:oneboundary}
设定理 \ref{thm:oneRHP} 的条件成立, $N$ 按 \eqref{eq:Nfactor} 定义, $n_{jk}=N(a)_{jk}$. 则物理场
\begin{equation}\label{eq:finalfields}
\begin{aligned}
\widetilde q&=-e^{-\ii\Omega t}(q_0+4bP_{12}),\\
\widetilde p&=\dfrac{e^{-\ii\Omega t}(N(a)R_\infty N(a)^\dagger)_{12}}{a^2+b^2},\qquad
\eta=\dfrac{(N(a)R_\infty N(a)^\dagger)_{11}}{a^2+b^2}
\end{aligned}
\end{equation}
在全体实 $x,t$ 上光滑, 满足原 HMB 系统 \eqref{eq:HMB}、Bloch 约束 \eqref{eq:C} 及两端同相边界 \eqref{eq:targetboundary}. 两个共振原像重构出同一个 Hermitian 迹零介质矩阵. 介质分量为
\begin{align}
\eta&=\dfrac{\eta_0}{a^2+b^2}\left[|n_{11}|^2-|n_{12}|^2
+\dfrac{\ii q_0}{a}(n_{11}\conj{n_{12}}-n_{12}\conj{n_{11}})\right],\label{eq:etacomponent}\\
e^{\ii\Omega t}\widetilde p&=\dfrac{\eta_0}{a^2+b^2}\left[n_{11}\conj{n_{21}}-n_{12}\conj{n_{22}}
+\dfrac{\ii q_0}{a}(n_{11}\conj{n_{22}}-n_{12}\conj{n_{21}})\right].\label{eq:pcomponent}
\end{align}
固定 $t$ 时, 三个场的端态误差为 $O(e^{-2\kappa|x|})$.
\end{theorem}
\begin{proof}
在无穷远展开 \eqref{eq:explicitM}, 得
\[
M^{(1)}=H(2\ii bI-4\ii bP-\ii q_0\sigma_1-2\ii\kappa\sg)H^{-1}.
\]
于是规范光场 $q=\ii M^{(1)}_{12}=-q_0-4bP_{12}$. 在实共振点, 由 $Y\lambda\sg Y^{-1}=\ii U_\infty$ 及 $N(a)^{-1}=\dfrac{N(a)^\dagger}{a^2+b^2}$,
\begin{equation}\label{eq:direct-medium}
M(z_a)\dfrac{\eta_0\lambda(z_a)}{a}\sg M(z_a)^{-1}
=\dfrac{HN(a)R_\infty N(a)^\dagger H^{-1}}{a^2+b^2}=:R.
\end{equation}
右端与谱叶无关. $N(a)N(a)^\dagger=(a^2+b^2)I$, 因而该酉共轭保持 Hermitian 性、迹零及 $R_\infty$ 的特征值. 特别地,
$R^2=\eta_0^2\left(1+\dfrac{q_0^2}{a^2}\right)I$, 给出 Bloch 约束. 结合 \eqref{eq:original-fields} 得到 \eqref{eq:finalfields}; 矩阵相乘给出两条分量式.

当 $x\to\pm\infty$ 时, $|T|$ 分别趋于无穷和零, 因而
\begin{equation}\label{eq:Pends}
P_+=\dfrac{1}{1+h_0^2}\begin{pmatrix}1&-h_0\\-h_0&h_0^2\end{pmatrix},\qquad
P_-=\dfrac{1}{1+h_0^2}\begin{pmatrix}h_0^2&-h_0\\-h_0&1\end{pmatrix}.
\end{equation}
这给出 $q\to q_0$. 将端态代入 $N$, 有
\[
N_+Y=Y_{-q_0}\diag(\lambda-\ii\kappa,\lambda+\ii\kappa),\qquad
N_-Y=Y_{-q_0}\diag(\lambda+\ii\kappa,\lambda-\ii\kappa),
\]
其中 $Y_{-q_0}=I+\dfrac{\ii q_0}{z}\sigma_1$. 利用 $HY_{-q_0}H^{-1}=Y$, 得
\begin{equation}\label{eq:direct-M-ends}
M\longrightarrow Y\quad(x\to-\infty),\qquad
M\longrightarrow Y\diag(s_{22},s_{22}^{-1})\quad(x\to+\infty).
\end{equation}
对角因子与 $\sg$ 对易, 所以两端 $R\to R_\infty$. 由 $|T|=e^{2\kappa[x-\operatorname{Re}F(\ii b)t]}$, 端态误差及任意固定阶导数均按 $e^{-2\kappa|x|}$ 衰减, 对紧 $t$ 区间一致. $\Delta_->0$ 与 $a^2+b^2>0$ 保证全局光滑.

对避开特殊点的 $z$ 定义
\begin{equation}\label{eq:inverse-log-lax}
U=(M_x-\ii\lambda M\sg)M^{-1},\qquad
V=(M_t+\ii\lambda F M\sg)M^{-1}.
\end{equation}
它们是 $\phi=Me^{\ii\theta\sg}$ 的两个对数导数. 换叶时 $\lambda$ 和 $\theta$ 变号, $\sigma_1$ 与 $\sg$ 反对易, 故两矩阵在 $z$ 和 $h(z)$ 取值相同, 是 $k$ 的单值有理矩阵.

在 $\zeta_j$ 附近, $\phi$ 的列留数为 $\Res\phi_2=c_j\phi_1(\zeta_j)$. 常数 $c_j$ 与 $x,t$ 无关, 故
\[
\phi=H_j(z)\left(I+\dfrac{c_je_1e_2^T}{z-\zeta_j}\right),\qquad
\det H_j(\zeta_j)=\gamma(\zeta_j)\ne0.
\]
右因子不依赖 $x,t$, 因而两个对数导数的离散极点可去; 共轭极点同理.

在分支点 $z_*=\pm\ii q_0$, 换叶约化使 $M(z_*)$ 消灭 $C(z_*)$ 的固定 $-1$ 特征空间. 该核不依赖 $x,t$, 因此 $M_x(z_*)$ 和 $M_t(z_*)$ 也消灭它. 又 $\lambda(z_*)=0$, $F$ 在此有限, \eqref{eq:inverse-log-lax} 的两个分子均消灭同一核. $\det M=\gamma$ 在 $z_*$ 有一阶零点, 而 $\operatorname{adj}M(z_*)$ 的像正是该核, 故两个商的分子抵消行列式的零点. 这证明分支点可去. 轮廓交点处 $\gamma\ne0$, 没有逆矩阵奇性. 在共振点, $U$ 解析, $V$ 的唯一留数为
\[
\Res_{k=a}V=\dfrac{\ii\eta_0\lambda(z_a)}{a}M(z_a)\sg M(z_a)^{-1}=\ii R.
\]

在第一叶无穷远, $M=I+O(k^{-1})$, $M_x,M_t=O(k^{-1})$. 由 \eqref{eq:reconq},
\[
U=-\ii k\sg+Q+O(k^{-1}).
\]
有限点均可去, Liouville 定理给出 $U=-\ii k\sg+Q$. 恒等式
\[
V=-FU+(M_t+FM_x)M^{-1}
\]
的第二项至多一次增长, 因而
\begin{equation}\label{eq:inverse-V-form}
V=4\ii\beta k^3\sg+k^2(-\ii\alpha\sg-4\beta Q)
+kW_1+W_0+\dfrac{\ii R}{k-a}.
\end{equation}
由 $\det\phi=\gamma$ 与 $x,t$ 无关, 两矩阵均迹零. 实谱上的共轭约化给出 $M^\dagger M=\gamma I$, 因此 $U,V$ 在实 $k\ne a$ 上反 Hermitian, 各系数具有相应约化.

作为同一 $\phi$ 的对数导数, 两矩阵恒满足零曲率关系. 比较 $k^2$ 项的非对角元及 $k$ 项的对角元, 得
\[
\begin{aligned}
(W_1)_{12}&=\alpha q-2\ii\beta q_x,\\
\partial_x(W_1)_{11}&=-2\ii\beta\partial_x|q|^2.
\end{aligned}
\]
故 $(W_1)_{11}=-2\ii\beta|q|^2+\chi_1(t)$. 由 \eqref{eq:direct-M-ends}, 两端 $V\to-FU_\infty$, 从而 $\chi_1(t)=0$. 再比较 $k$ 项的非对角元及常数项的对角元, 有
\[
\begin{aligned}
(W_0)_{12}&=2\beta|q|^2q+\dfrac{\ii\alpha}{2}q_x+\beta q_{xx},\\
\partial_x(W_0)_{11}&=\partial_x\left[\dfrac{\ii\alpha}{2}|q|^2-\beta(q\conj{q_x}-q_x\conj{q})\right].
\end{aligned}
\]
同一端态固定后一积分常数为 $\dfrac{\ii\Omega}{2}$. 于是 $W_1=V_1(q)$、$W_0=V_0(q)+\dfrac{\ii\Omega}{2}\sg$, 其余分量由迹零和约化确定, 得到完整规范 Lax 对 \eqref{eq:gaugelax}.

零曲率的共振留数和常数非对角元分别给出
\[
R_x=[-\ii a\sg+Q,R],\qquad
q_t=\ii\alpha\left(\dfrac{q_{xx}}{2}+|q|^2q\right)
+\beta(q_{xxx}+6|q|^2q_x)+2p+\ii\Omega q.
\]
前式为两条介质方程; 恢复 \eqref{eq:original-fields} 即得原 HMB 系统.

最后, 在普通连续谱上令
\[
\phi_-=Me^{\ii\theta\sg},\qquad
\phi_+=M\diag(s_{22}^{-1},s_{22})e^{\ii\theta\sg}.
\]
已恢复的 Lax 对及 \eqref{eq:direct-M-ends} 表明两者具有相应端点归一化, 由 Jost 唯一性即为重构场的 Jost 解. 所以 $S=\diag(s_{22}^{-1},s_{22})$, 反射为零, 透射函数为 \eqref{eq:onezeros}; 外部零点的列关系给出 $b_1=c_1s_{22}'(\zeta_1)=1$.
\end{proof}

写 $s=\xi+\ii\nu$, 光场公式中的投影元还可表示为
\begin{equation}\label{eq:hyperbolic-projector}
P_{12}=\dfrac{-2h_0\cosh(2\xi)+(1+h_0^2)\cos(2\nu)+\ii(1-h_0^2)\sin(2\nu)}
{2(1+h_0^2)\cosh(2\xi)-4h_0\cos(2\nu)}.
\end{equation}
分母不小于 $2(1-h_0)^2$, 局域包络与内部振荡分别由 $\xi$ 和 $\nu$ 控制.
